\section{How the System Works}
\label{sec:system}

\subsection{Data, Shards, and Summaries}

% CITE-PLAIN: paper/sections/setup.tex:36 (nguyen2016msmarco), first part.
We use \RNumCorpusTotalPassages{} passages and \RNumPlainTestQueries{} test queries from the MS MARCO data set \cite{nguyen2016msmarco}.
% CITE-PLAIN: paper/sections/setup.tex:36 (nguyen2016msmarco), second part; "real" and "Bing" dropped.
Its questions come from anonymized user queries, and its passages come from web documents \cite{nguyen2016msmarco}.
% CITE-PLAIN: paper/sections/setup.tex:36 (reimers2019sbert); "six-layer" dropped.
A language model called MiniLM, from the Sentence-Transformers software package \cite{reimers2019sbert}, turns each passage and query into a vector of \RNumVectorDimFloats{} numbers, or \RNumPlainVectorBytes{} bytes.
We compare two vectors by multiplying their matching numbers and adding the products; a higher score means closer meaning.

The $k$-means method sorts vectors into $k$ groups of similar vectors and finds the center, or average point, of each group.
% CITE-PLAIN: paper/sections/system.tex:4 (douze2024faiss); Faiss glossed from the LEDGER quote.
At setup, $k$-means from Faiss, a software package for searching and grouping vectors \cite{douze2024faiss}, splits the starting passages into the \RNumNumShards{} shards.
We call the center of a shard's group its starting center.
Each shard then runs $k$-means with \RNumCentersPerSummary{} groups on its own passages, and its summary is the list of their centers, \RNumSummaryBytes{} bytes.

To answer a query, the coordinator scores each shard by the summary center closest to the query, and it searches the two best-scoring shards.
% CITE-PLAIN: paper/sections/system.tex:4 (malkov2020hnsw); "proximity graph" glossed, the method name dropped.
Each shard stores its vectors in a search graph that links each vector to nearby ones \cite{malkov2020hnsw}.
A search hops along these links toward the query, and it can miss a few close passages.
Each searched shard returns the IDs and scores of the \RNumPlainTopK{} closest passages it finds, and the coordinator keeps the best \RNumPlainTopK{}.
New passages join a shard without moving old ones, and the coordinator keeps the shard's old summary until the shard sends a fresh one.

Each time step draws its queries without repeats from the \RNumPlainTestQueries{} test queries, \RNumPlainQueriesDefault{} by default.
Under starting-center placement, described below, we also test \RNumPlainQueriesLow{} and \RNumPlainQueriesHigh{} queries per time step.

We test two streams, that is, two orders in which new passages arrive.
In the mixed-topic stream, half the passages build the shards, and the other half arrive in random order, \RNumSteadyEpochInsertVecs{} per time step.
In the new-topic stream, we first sort all passages into \RNumPlainTopicClusters{} clusters of similar passages with $k$-means, so each cluster is roughly one topic.
Half of the clusters, picked at random, build the shards.
The other half arrive one cluster after another, about \RNumPlainNewTopicStepVecs{} passages per time step on average, bringing topics the shards have not seen.
% CITE-PLAIN: paper/sections/setup.tex:36 (simhadri2024bigann); "into 64 clusters" dropped.
This design is modeled on the streaming test of the Big-ANN NeurIPS'23 vector-search competition, which inserts points cluster by cluster \cite{simhadri2024bigann}.

Passages never move once stored, and we test three ways to place new ones.
In starting-center placement, each new passage goes to the shard with the closest starting center.
In one-shard placement, all of a time step's new passages go to one shard, and the next time step uses the next shard.
Every refresh rule then refreshes that shard, because rules serve shards with new passages first (Section~\ref{sec:rules}).
In eight-shard placement, a group of eight shards is open in each time step, and the group moves on by eight shards every step.
Each new passage goes to the open shard with the closest starting center, so each step changes five to eight shards.
In these last two placements, the coordinator picks shards only by their summaries and does not know which shards are open.

\begin{table}[!t]
\centering
\caption{What Each Option Sends and What a Budget of 4 Buys per Time Step}
\label{tab:options}
\footnotesize
\setlength{\tabcolsep}{3pt}
\begin{tabular}{@{}l>{\raggedright\arraybackslash}p{4.2cm}r>{\raggedright\arraybackslash}p{1.55cm}@{}}
\toprule
Item & What is sent & Bytes & Budget of 4 buys \\
\midrule
Refresh & One shard's summary (\RNumCentersPerSummary{} centers) & \RNumSummaryBytes{} & \RNumFourRulesSummariesBFour{} summaries \\
Extra search & A query to one more shard, and IDs and scores of its \RNumPlainTopK{} best matches & \RNumPlainSearchBytes{} & Up to \RNumExtraSearchesBFour{} searches \\
Push & A new passage and its ID number & \RNumPlainPushBytes{} & Up to \RNumPushedVecsBFourPerEpoch{} passages \\
\midrule
Probe & Its scores, cutoffs, and counts in one time step (Section~\ref{sec:rules}) & \RNumRNOverheadEpochBytes{} & Room for \RNumRNSummariesBFour{} summaries \\
\bottomrule
\end{tabular}
\end{table}

\subsection{The Budget and Its Three Uses}

The budget counts the network bytes each time step may use, in units of one summary.
We test budgets of two, four, and eight summaries; four is the default and allows \RNumBudgetBFourEpochBytes{} bytes (Table~\ref{tab:options}).
To refresh, a shard rebuilds its summary over all its passages and sends it.
Refreshing all \RNumNumShards{} shards every time step would cost \RNumAlwaysFreshEpochBytes{} bytes, eight times the default budget.

Extra search spends the budget on the queries where the coordinator is least sure.
These are the queries whose second- and third-ranked shards have the closest scores, so either could hold the best matches.
For such a query, the coordinator also searches the third-ranked shard.
The planned version picks these queries after seeing all of the time step's queries, which a system answering each query on arrival could not do.
A budget of four buys up to \RNumExtraSearchesBFour{} extra searches, so at \RNumPlainQueriesLow{} queries per step it also asks a fourth shard for some queries.
The live version decides as each query arrives and asks at most one extra shard per query.
It asks when the score gap is below a cutoff: the gap that would have spent the budget exactly in the previous time step.
With push, shards send up to the budget's worth of their newest passages each time step.
The coordinator keeps them all and checks each query against them.

We also compare against two reference methods that ignore the budget.
Never refresh keeps the first summaries, and refresh all sends all \RNumNumShards{} fresh summaries every time step.
Refresh all is a reference, not the best that refresh can do.
Its rebuilt summaries still have only \RNumCentersPerSummary{} centers, so they describe a shard roughly.
Searching all \RNumNumShards{} shards finds more than \RNumPlainAllShardsMinPct{} of recent matches, so almost all misses come from choosing shards.

\subsection{Five Refresh Rules}
\label{sec:rules}

A refresh rule picks which shards refresh in each time step.
Every rule serves first the shards with new passages since their last refresh, putting any that waited too long ahead.
The waiting limit is twice the time a rule needs to refresh every shard once.
That is \RNumPlainCapBFour{} time steps for a rule that refreshes four summaries per step.
Any budget left over re-sends the summaries of unchanged shards.
Rules that score shards break ties by the lowest shard number.

Rotation lets shards take turns.
Most-searched favors the shards that past test queries were sent to most often, and largest-shard favors the shards with the most passages.
Their scores do not reset after a refresh, so they keep picking the same shards until the waiting limit steps in.
Most-changed scores a shard by the share of its passages that are new.
It multiplies that share by how far the new passages lie, on average, from the nearest center in the shard's last summary.
Each shard computes this score itself, so a real system would send one \RNumScoreBytesPerShard{}-byte score per shard, or \RNumDCScoreEpochBytes{} bytes per time step.
Our main runs did not charge these bytes; the follow-up in Section~\ref{sec:follow} does.

The sample-query probe is our own rule and the most complex.
% CITE-PLAIN: paper/sections/system.tex:61 (pandey2008crawlorder, olston2002besteffort); the two ideas glossed as in Section II.
It adapts two ideas from Section~\ref{sec:related}: ranking by how often items reach top results \cite{pandey2008crawlorder}, and letting each source rank its own changes \cite{olston2002besteffort}.
The probe uses \RNumSampleQueries{} sample queries: practice questions from a separate part of the data set, not the test queries.
The coordinator sends them to every shard once at setup.
Each time step then has three rounds of messages.
First, for each sample query, the coordinator asks the two shards that its summaries rank highest for the scores of their \RNumPlainTopK{} best matches.
Second, it sends every shard two cutoffs for each sample query.
One is the summary score a shard would need to rank in the top two.
The other is the score of the query's \RNumPlainTopK{}th-best match so far.
Third, each shard with new passages rebuilds its summary locally.
For each sample query, it checks whether the rebuilt summary would lift it into the top two.
For those queries, it counts its new passages that score above the \RNumPlainTopK{}th-best cutoff and sends that count.
Among shards with new passages, those over the waiting limit refresh first.
The rest go in order of their counts, with ties going to the lowest shard number.

These rounds cost \RNumRNOverheadEpochBytes{} bytes per time step.
Each budget is a whole number of summaries, and refresh can send only whole summaries.
So any extra charged cost, however small, uses up the room for one summary.
At a budget of four, the probe therefore refreshes three summaries where the other rules refresh four, and its waiting limit is \RNumPlainCapRNBFour{} time steps.

\subsection{Costs We Do Not Count}

The budget pays only for the messages in Table~\ref{tab:options}, and six other costs are left out.
The first is sending new passages to the shards, which is the same for every method.
The second is the probe's setup broadcast of \RNumRNSetupBroadcastBytes{} bytes, its small per-step lists of which sample queries to answer, and its shard work.
Every changed shard rebuilds a summary just to compute its count, and the shards run \RNumRNSampleSearchesPerEpoch{} sample-query searches per time step.
Spread over a run, the broadcast would add \RNumRNSetupAmortEpochBytes{} bytes per step.
That is more than a budget of two, and at four it would leave no room for any summary.
The third is most-changed's score bytes, in the main runs.
The fourth is a free flag that tells the coordinator which shards have new passages.
The flag says only that a shard has new passages, not what they are about.
Every refresh rule uses it to choose shards to refresh, but queries are still sent by summaries alone.
Largest-shard also gets the shard sizes for free.
Charging even a few bytes for this flag would cost every refresh rule one whole summary per time step.
The fifth is the added work of extra search: at the default query rate, shards run \RNumESLoadPctBTwo{} to \RNumESLoadPctBEight{} more searches across budgets.
At a budget of four and \RNumPlainQueriesLow{} queries per step, they run up to \RNumESLoadPctQFifty{} more.
Asking a third shard also means contacting three shards at once or waiting for a second exchange of messages, and neither is counted.
The sixth is push's memory and scanning, up to \RNumPVMaxVecsBFour{} kept passages at a budget of four and \RNumPVMaxVecsBEight{} at eight.
Scanning them for each query takes \RNumPlainPVScanRatioBFour{} to \RNumPlainPVScanRatioBEight{} times as many vector comparisons as scoring all summaries.
Keeping only the last two time steps' \RNumPVFifoVecsBFour{} would give the same recent recall at a budget of four.
Push also assumes the shards know which passages arrived last overall, since it sends the newest ones.
Leaving out the probe's extra costs, most-changed's score bytes, and the free flag favors the refresh rules.
Leaving out the extra search work and push's memory and scanning favors extra search and push.
